% GENERATED FILE. Edit canonical lessons, then run npm run generate:books.
% canonical-source-hash: fnv1a64:cbc998964a78bcdb
% canonical-lessons: MW-C161-mhai, MW-C161-kiyo-liyo, MW-C161-kadhno, MW-W49-ddha, MW-R161-what-i-did

\chapter{What I Did}
\label{ch:mw-a2-what-i-did}
\hlchaptermodality{161}

\begin{chapteropening}
\textbf{By the end of this chapter:} \emph{I can say what I ate, drank, did and took, with the ending that agrees with the thing, and write the letter \mw{ढ} from \mw{काढणो}.}

Complete MW-R161-what-i-did: say what you ate, drank, did, took and drew, with \mw{म्हैं}, and write \mw{ढ}.
\end{chapteropening}

\section[\texorpdfstring{\mw{म्हैं} \mw{रोटी} \mw{खाई।} (mhaĩ roṭī khāī.)}{म्हैं रोटी खाई। (mhaĩ roṭī khāī.)}]{\mw{म्हैं} \mw{रोटी} \mw{खाई} — I ate}

\label{lesson:MW-C161-mhai}



\begin{quote}
Say \emph{I eat roti} and \emph{I drink water}. Then \emph{mother went home}.
\end{quote}

\begin{grammarlens}[title={Grammar Lens: \mw{म्हैं}}]
Eating, drinking and seeing are done \emph{to} something. In their past, \emph{I} is \textbf{\mw{म्हैं}}, not \textbf{\mw{हूं}}:

\begin{quote}
\textbf{\mw{म्हैं} \mw{रोटी} \mw{खाई।}}
\end{quote}

\begin{quote}
\emph{mhaĩ roṭī khāī.} — I ate a roti.
\end{quote}

Old Marwari already had \emph{mẽ}, \textquotedblleft{}I\textquotedblright{}, beside \emph{hũ} (CDIAL 9691, 992).
\end{grammarlens}

\begin{grammarlens}[title={Grammar Lens: the thing decides}]
The ending follows the thing, not you. \textbf{\mw{रोटी}} is feminine, so \textbf{\mw{खाई}}; \textbf{\mw{पाणी}} is masculine, so \textbf{\mw{पीयो}}:

\begin{quote}
\textbf{\mw{म्हैं} \mw{पाणी} \mw{पीयो।}}
\end{quote}

\begin{quote}
\emph{mhaĩ pāṇī pīyo.} — I drank water.
\end{quote}

A man and a woman say the same: the roti decides, and so does the water. This past, with \textbf{\mw{म्हैं}} and an ending that follows the thing, is central Marwari as this book writes it; other districts differ.
\end{grammarlens}

\subsection*{Your turn}
Say these aloud:

\begin{itemize}
\raggedright
  \item \textbf{\mw{म्हैं} \mw{दूध} \mw{पीयो।}} — \emph{mhaĩ dūdh pīyo.} — I drank milk.
  \item \textbf{\mw{म्हैं} \mw{गाय} \mw{देखी।}} — \emph{mhaĩ gāy dekhī.} — I saw the cow.
  \item \textbf{\mw{म्हैं} \mw{घर} \mw{देख्यो।}} — \emph{mhaĩ ghar dekhyo.} — I saw the house.
\end{itemize}

\begin{tcolorbox}[breakable,colback=teal!4,colframe=teal!35!black,title={Before you move on}]
\emph{I}, in the past of eating? (\textbf{\mw{म्हैं}}.) Who decides \textbf{\mw{खाई}} or \textbf{\mw{खायो}}? (The thing eaten.)
\end{tcolorbox}
\section[\texorpdfstring{\mw{कियो} · \mw{लियो} (kiyo · liyo)}{कियो · लियो (kiyo · liyo)}]{\mw{कियो}, \mw{लियो} — did, took}

\label{lesson:MW-C161-kiyo-liyo}



\begin{quote}
Say \emph{I ate a roti}. Then \emph{to do} and \emph{to take}: \textbf{\mw{करणो}}, \textbf{\mw{लेणो}}.
\end{quote}

\begin{grammarlens}[title={Grammar Lens: two short pasts}]
\textbf{\mw{करणो}} and \textbf{\mw{लेणो}} have short pasts: \textbf{\mw{कियो}} and \textbf{\mw{लियो}}.

\begin{quote}
\textbf{\mw{म्हैं} \mw{पाणी} \mw{गरम} \mw{कियो।}}
\end{quote}

\begin{quote}
\emph{mhaĩ pāṇī garam kiyo.} — I heated the water.
\end{quote}

\begin{quote}
\textbf{\mw{म्हैं} \mw{दूध} \mw{लियो।}}
\end{quote}

\begin{quote}
\emph{mhaĩ dūdh liyo.} — I took milk.
\end{quote}

For a feminine thing, \textbf{\mw{लियो}} turns into \textbf{\mw{ली}}:

\begin{quote}
\textbf{\mw{म्हैं} \mw{रोटी} \mw{ली।}}
\end{quote}

\begin{quote}
\emph{mhaĩ roṭī lī.} — I took a roti.
\end{quote}

CDIAL lists Old Marwari \emph{kiyo}, \textquotedblleft{}made\textquotedblright{}, and \emph{liyo}, \textquotedblleft{}takes\textquotedblright{} (CDIAL 2814, 10948). This book pairs \textbf{\mw{कियो}} with a masculine thing only. \textbf{\mw{ली}} is central Marwari as this book writes it; other districts differ.
\end{grammarlens}

\subsection*{Your turn}
The thing changes the ending; the speaker does not:

\noindent
\begin{tabularx}{\linewidth}{@{}>{\raggedright\arraybackslash}X>{\raggedright\arraybackslash}X>{\raggedright\arraybackslash}X@{}}
\toprule
\textbf{who speaks} & \textbf{the thing} & \textbf{the verb} \\
\midrule
a man & \textbf{\mw{रोटी}} & \textbf{\mw{ली}} \\
a woman & \textbf{\mw{रोटी}} & \textbf{\mw{ली}} \\
a man & \textbf{\mw{दूध}} & \textbf{\mw{लियो}} \\
a woman & \textbf{\mw{दूध}} & \textbf{\mw{लियो}} \\
\bottomrule
\end{tabularx}

Say these aloud:

\begin{itemize}
\raggedright
  \item \textbf{\mw{म्हैं} \mw{रोटी} \mw{ली।}} — \emph{mhaĩ roṭī lī.} — I took a roti.
  \item \textbf{\mw{म्हैं} \mw{दूध} \mw{लियो।}} — \emph{mhaĩ dūdh liyo.} — I took milk.
  \item \textbf{\mw{म्हैं} \mw{पाणी} \mw{गरम} \mw{कियो।}} — \emph{mhaĩ pāṇī garam kiyo.} — I heated the water.
\end{itemize}

\begin{tcolorbox}[breakable,colback=teal!4,colframe=teal!35!black,title={Before you move on}]
\emph{Did}, of water? (\textbf{\mw{कियो}}.) \emph{Took}, of milk? (\textbf{\mw{लियो}}.) \emph{Took}, of a roti? (\textbf{\mw{ली}}.)
\end{tcolorbox}
\section[\texorpdfstring{\mw{काढणो} (kāḍhṇo)}{काढणो (kāḍhṇo)}]{\mw{काढणो} (kāḍhṇo) — to take out}

\label{lesson:MW-C161-kadhno}



\begin{quote}
Before the new one: say \emph{to pull} and \emph{to lift}. (\textbf{\mw{खींचणो}}, \textbf{\mw{उठावणो}}.) Then \emph{water}.
\end{quote}

\subsection*{You'll want to know: \mw{काढणो}}
\textbf{\mw{काढणो}} — \emph{kāḍhṇo} — \textquotedblleft{}to take out; to draw (water)\textquotedblright{}.

\begin{quote}
\textbf{\mw{म्हैं} \mw{पाणी} \mw{काढ्यो।}}
\end{quote}

\begin{quote}
\emph{mhaĩ pāṇī kāḍhyo.} — I drew water.
\end{quote}

\begin{quote}
\textbf{\mw{वो} \mw{रोज} \mw{पाणी} \mw{काढै} \mw{है।}}
\end{quote}

\begin{quote}
\emph{vo roj pāṇī kāḍhai hai.} — He draws water every day.
\end{quote}

It holds a letter you have not written yet: \emph{ḍh}, a \textbf{\mw{ड}} with a breath after it. The next lesson puts it in your hand.

\begin{cousinweb}
\textbf{\mw{काढणो}} (\emph{kāḍhṇo}) is inherited from a reconstructed verb, \emph{kaḍḍhati}, \textquotedblleft{}pulls, draws\textquotedblright{}. CDIAL lists Old Marwari \emph{kāḍhaï}, \textquotedblleft{}to pull out\textquotedblright{} (CDIAL 2660).
\end{cousinweb}

\begin{grammarlens}[title={What you've built}]
Its stem \textbf{\mw{काढ}} ends in a consonant: \textbf{\mw{काढूं} \mw{हूं}}, \textbf{\mw{काढै} \mw{है}}, and in the past, with \textbf{\mw{म्हैं}}, \textbf{\mw{काढ्यो}}.
\end{grammarlens}

\subsection*{Your turn}
\begin{itemize}
\raggedright
  \item \emph{Say it:} \emph{kāḍhṇo}
  \item \emph{Say it:} \textbf{\mw{म्हैं} \mw{पाणी} \mw{काढ्यो।}} — \emph{mhaĩ pāṇī kāḍhyo.} — I drew water.
  \item \emph{Recall:} say \emph{to pull}, then \emph{kāḍhṇo} again
\end{itemize}

\begin{tcolorbox}[breakable,colback=teal!4,colframe=teal!35!black,title={Before you move on}]
What does \textbf{\mw{काढणो}} mean? (\textquotedblleft{}To take out; to draw water\textquotedblright{}.) Say it once more.
\end{tcolorbox}
\section[\texorpdfstring{\mw{ढ} (ḍha)}{ढ (ḍha)}]{\mw{ढ} (ḍha) — the letter inside \mw{काढणो}}

\label{lesson:MW-W49-ddha}



\begin{quote}
Before the new one: write \textbf{\mw{ड}} once, and name it.

Now say \textbf{\mw{काढणो}} (\emph{kāḍhṇo}), \textquotedblleft{}to take out\textquotedblright{}. Its middle is a \textbf{\mw{ड}} with a breath after it: \emph{ḍh}. This lesson puts that letter in your hand.
\end{quote}

\subsection*{Script you'll notice: \mw{ढ}}
\textbf{\mw{ढ}} — \emph{ḍha}, the breathed \emph{ḍ}.

It sits inside \textbf{\mw{काढणो}}. It stands beside \textbf{\mw{ड}} and \textbf{\mw{ड़}} as \textbf{\mw{ध}} stands beside \textbf{\mw{द}}. Its body is one stroke: a stem that runs into a broad outer bowl and a closed loop inside it, then the headline.

\subsection*{Writing: \mw{ढ}}
\hlblockfigure{\detokenize{figures/MW-W49-ddha-filmstrip.pdf}}{How \mw{ढ} is written, stroke by stroke}

\begin{itemize}
\raggedright
  \item \textbf{1.} start at the stem's headline junction, descend, turn left through the upper shoulder, sweep counterclockwise around the broad outer bowl, and continue counterclockwise around the closed inner loop without lifting
  \item \textbf{2.} lift and draw the top shirorekhā left-to-right
\end{itemize}

\textbf{Pen lifts: 1.}

\begin{quote}
verified two-stroke teaching form fitted to the bundled printed outline
\end{quote}

\begin{quote}
Stroke order is one attested teaching order; hands and teachers vary. Source: Opiaterein, ‘Deva-\mw{ढ}-order.gif’, strokes 1–2, Wikimedia Commons, 10 May 2009.
\end{quote}

\subsection*{Your turn}
\begin{itemize}
\raggedright
  \item \emph{Look:} at \textbf{\mw{काढणो}} and point at \mw{ढ} inside it
  \item \emph{Trace it:} \mw{ढ} three times, saying \emph{ḍh} as you finish each one
  \item \emph{Write it:} \mw{ड}, then \mw{ढ}, from memory
  \item \emph{Read it:} \textbf{\mw{काढणो}} aloud once more
\end{itemize}

\begin{tcolorbox}[breakable,colback=teal!4,colframe=teal!35!black,title={Before you move on}]
Which letter is this — \mw{ढ}? (\textbf{ḍha}, the breathed \emph{ḍ}.) Which word did it come from? (\textbf{\mw{काढणो}}, \emph{kāḍhṇo}, \textquotedblleft{}to take out\textquotedblright{}.)
\end{tcolorbox}
\section[(dialogue)]{What I did — a review}

\label{lesson:MW-R161-what-i-did}



\begin{quote}
Say \emph{I ate a roti}.
\end{quote}

\subsection*{Your turn}
\begin{itemize}
\raggedright
  \item \emph{Say it:} \emph{I ate a roti}; \emph{I drank water}; \emph{I drank milk}
  \item \emph{Say it:} \emph{I saw the cow}; \emph{I saw the house}
  \item \emph{Say it:} \emph{I heated the water}; \emph{I heated the milk}
  \item \emph{Say it:} \emph{I took a roti}; \emph{I took milk}
  \item \emph{Say it:} \emph{I drew water}; \emph{he draws water every day}
  \item \emph{Read it:} \textbf{\mw{काढणो}}, and name the letter after \textbf{\mw{का}}
\end{itemize}

\begin{tcolorbox}[breakable,colback=teal!4,colframe=teal!35!black,title={Before you move on}]
\emph{I ate a roti}? (\textbf{\mw{म्हैं} \mw{रोटी} \mw{खाई।}}) \emph{I took milk}? (\textbf{\mw{म्हैं} \mw{दूध} \mw{लियो।}}) \emph{I drew water}? (\textbf{\mw{म्हैं} \mw{पाणी} \mw{काढ्यो।}})
\end{tcolorbox}
